% !PRINTED 349
\weberpage{349}

%% [running head, not body text -- commented out 2026-08-07] Page furniture (running
%% head, folio, signature) is excluded from this edition by standing policy.
%% \noindent\S\ 95.\hfill Equivalent forms and ideals in the quadratic field.\hfill 349

Now let $\varphi = (a, b, c)$, $\varphi' = (a', b', c')$ be, conversely, two
\gesperrt{equivalent forms with positive first coefficients} $a, a'$, and
\weq{6}{\omega = -\frac{b+\sqrt{\Delta}}{2a}, \quad \omega' = -\frac{b'+\sqrt{\Delta}}{2a'}\,.}
Then there is a linear substitution $\begin{pmatrix} \alpha, & \beta \\ \gamma, & \delta \end{pmatrix}$ with determinant $+1$ by which $(a, b, c)$ passes into $(a', b', c')$.

We then have
\weq{7}{
\begin{aligned}
a' &= a\alpha^2 + b\alpha\gamma + c\gamma^2,\\
b' &= 2a\alpha\beta + b(\alpha\delta + \beta\gamma) + c\gamma\delta,\\
c' &= a\beta^2 + b\beta\delta + c\delta^2,
\end{aligned}
}
\weq{8}{
\omega = \frac{\alpha\omega' + \beta}{\gamma\omega' + \delta}, \quad \omega' = \frac{\delta\omega - \beta}{-\gamma\omega + \alpha}\,,
}
and a short calculation shows that $\omega$ and $\omega'$ are corresponding roots, i.\,e. roots with the same sign of $\sqrt{\Delta}$.

Accordingly, for the ideals $\mathfrak{a}, \mathfrak{a}'$ belonging to $\varphi$ and $\varphi'$ we obtain two basis forms:
\[
\begin{aligned}
\lambda &= a(x - \omega y),\\
\lambda' &= a'(x' - \omega'y'),
\end{aligned}
\]
and by (8) we have
\[
\lambda = \frac{a[(\delta x - \beta y) + (\gamma x - \alpha y)\omega']}{\gamma\omega' + \delta}\,.
\]

If we therefore set
\[
x' = \delta x - \beta y, \quad y' = -\gamma x + \alpha y,
\]
it follows that
\[
\lambda = \frac{a}{a'(\gamma\omega' + \delta)}\,\lambda',
\]
and since $a : a'(\gamma\omega' + \delta)$ is a number in $\Omega$, it follows that $\mathfrak{a}$ and $\mathfrak{a}'$ are equivalent.

We thus have once more, in agreement with the general theorem \S\ 88, 4. and at the same time in closer determination:

6. \gesperrt{If $(a, b, c)$ $(a', b', c')$ are two equivalent forms with positive first coefficients, and if $\omega, \omega'$ are corresponding roots of these forms, then the ideals $(a, a\omega)$ and $(a', a'\omega')$ are equivalent.}

For a negative discriminant the first coefficients of equivalent forms always have the same sign, and one considers only those forms in which this sign is positive.
%CONT>>
